\uextra{Apéndice}{Catálogo de Nodos del Lenguaje}
\label{apx:node-catalog}

\begingroup
\renewcommand{\TableBodyFontSize}{\scriptsize}

\begin{table}[H]
    \centering
    \caption{Nodos de Entrada}
    \begin{tabularx}{\textwidth}{ l X X X c }
        \hline
        \textbf{Nodo}    & \textbf{Entradas}      & \textbf{Proceso}                                                                    & \textbf{Salida}                                     & \textbf{Nivel de Abstracción} \\ \hline
        Leer             & Ninguna                & Pausa la ejecución para esperar que el usuario ingrese un dato por teclado.         & El dato ingresado por el usuario (Número o texto).  & 1 \\
        Número           & Ninguna (configurable) & Mantiene un valor numérico fijo                                                     & El número configurado                               & 1 \\
        Texto            & Ninguna (configurable) & Mantiene un valor de texto fijo                                                     & El texto configurado                                & 1 \\
        Valor lógico     & Ninguna (configurable) & Mantiene un valor booleano fijo (Verdadero o Falso)                                 & El valor lógico configurado                         & 1 \\
        Lista fija       & Ninguna (configurable) & Mantiene una lista de elementos fija, definida por el usuario.                      & La lista configurada.                               & 3 \\
        Valor Aleatorio  & Ninguna (configurable) & Genera un número al azar dentro del rango establecido                               & Número aleatorio generado                           & 2 \\
        Opción Aleatoria & Ninguna (configurable) & Genera un valor aleatorio entre un conjunto de opciones predefinidas por el usuario & Opción seleccionada aleatoriamente (Número o Texto) & 2 \\
        Generar pulsos   & Inicio, Fin, Paso      & Genera una secuencia de números y emite cada valor como pulso individual            & Pulsos con cada valor de la secuencia               & 1 \\
        Temporizador     & Valor (cualquier tipo) & Espera la cantidad de segundos indicada por el usuario y luego se activa            & Valor de entrada                                    & 2 \\
        \hline
    \end{tabularx}
\end{table}

\begin{table}[H]
    \centering
    \caption{Nodos de Salida}
    \begin{tabularx}{\textwidth}{ l X X X c }
        \hline
        \textbf{Nodo} & \textbf{Entradas}         & \textbf{Proceso}                                        & \textbf{Salida}     & \textbf{Nivel de Abstracción} \\ \hline
        Escribir      & Datos (cantidad dinámica) & Muestra los valores recibidos en la consola del usuario & Sin salida de datos & 1 \\ \hline
    \end{tabularx}
\end{table}

\begin{table}[H]
    \centering
    \caption{Nodos de Matemática}
    \begin{tabularx}{\textwidth}{ l X X X c }
        \hline
        \textbf{Nodo}        & \textbf{Entradas}                                             & \textbf{Proceso}                                                                    & \textbf{Salida} & \textbf{Nivel de Abstracción} \\ \hline
        Suma                 & Número A, Número B                                            & Calcula la suma de A más B.                                                         & Número          & 1-2 \\
        Resta                & Número A, Número B                                            & Resta el valor de B al valor de A.                                                  & Número          & 1-2 \\
        Multiplicación       & Número A, Número B                                            & Multiplica el valor de A por el valor de B.                                         & Número          & 1-2 \\
        División             & Dividendo, Divisor                                            & Divide el Dividendo entre el Divisor.                                               & Número          & 1-2 \\
        Resto                & Dividendo, Divisor                                            & Calcula el resto entre el Dividendo y el Divisor.                                   & Número          & 1-2 \\
        Potencia             & Base, Exponente                                               & Eleva la Base a la potencia del Exponente.                                          & Número          & 2 \\
        Raíz                 & Radicando, Índice                                             & Calcula la raíz n-ésima del radicando según el índice indicado.                     & Número          & 2 \\
        Logaritmo            & Valor, Base                                                   & Calcula el logaritmo del valor en la base seleccionada.                             & Número          & 2 \\
        Piso                 & Número                                                        & Redondea el número hacia abajo al entero más cercano.                               & Número          & 2 \\
        Techo                & Número                                                        & Redondea el número hacia arriba al entero más cercano.                              & Número          & 2 \\
        Redondear            & Número                                                        & Redondea el número al entero más cercano.                                           & Número          & 2 \\
        Editor de fórmulas   & Cantidad definida por la expresión introducida por el usuario & Se realiza la expresión definida por el usuario.                                    & Número          & 3 \\
        Operación Matemática & Número A, Número B                                            & Se realiza la operación matemática seleccionada por el usuario en la configuración. & Número          & 2 \\
        \hline
    \end{tabularx}
\end{table}

\begin{table}[H]
    \centering
    \caption{Nodos de Manejo de Texto}
    \begin{tabularx}{\textwidth}{ l X X X c }
        \hline
        \textbf{Nodo} & \textbf{Entradas}                                      & \textbf{Proceso}                                                               & \textbf{Salida} & \textbf{Nivel de Abstracción} \\ \hline
        Concatenar    & Texto A, Texto B                                       & Une el Texto B al final del Texto A.                                           & Texto           & 1 \\
        Reemplazar    & Texto original, Texto a reemplazar, Texto de reemplazo & Reemplaza las ocurrencias de un subtexto por el texto de reemplazo.            & Texto           & 1 \\
        Sustraer      & Texto original, Texto a eliminar                       & Elimina de un texto el subtexto especificado.                                  & Texto           & 1 \\
        Recortar      & Texto original, Inicio, Fin                            & Extrae del texto original el texto entre la posición Inicio y la posición Fin. & Texto           & 1 \\
        \hline
    \end{tabularx}
\end{table}

\begin{table}[H]
    \centering
    \caption{Nodos de Listas: Construcción y Combinación}
    \begin{tabularx}{\textwidth}{ l l X X c }
        \hline
        \textbf{Nodo}        & \textbf{Entradas} & \textbf{Proceso}                                                                        & \textbf{Salida}       & \textbf{Nivel de Abstracción} \\ \hline
        Generar lista        & Inicio, Fin, Paso & Genera una lista de números desde Inicio hasta Fin con el Paso indicado                 & Lista numérica        & 1 \\
        Unir listas          & Lista A, Lista B  & Empaqueta los elementos de las dos listas en una lista de tuplas                        & Lista empaquetada     & 2 \\
        Separar listas       & Lista de tuplas   & Transforma una lista de tuplas en dos listas separadas                                  & Lista A, Lista B      & 2 \\
        Separar              & Número o Texto    & Convierte un número o texto en una lista de sus elementos                               & Lista                 & 2 \\
        \hline
    \end{tabularx}
\end{table}

\begin{table}[H]
    \centering
    \caption{Nodos de Listas: Transformación}
    \begin{tabularx}{\textwidth}{ l l X X c }
        \hline
        \textbf{Nodo}      & \textbf{Entradas} & \textbf{Proceso}                                                                        & \textbf{Salida}    & \textbf{Nivel de Abstracción} \\ \hline
        Mapear             & Lista             & Aplica una función a cada elemento de la lista y devuelve la lista transformada         & Lista transformada & 2 \\
        Filtrar            & Lista             & Conserva los elementos de la lista para los que una condición configurada sea verdadera & Lista filtrada     & 2 \\
        Ordenar            & Lista             & Ordena una lista según la configuración                                                 & Lista              & 3 \\
        Voltear            & Lista             & Invierte el orden de los elementos de la lista                                          & Lista              & 2 \\
        Posiciones pares     & Lista             & Devuelve los elementos ubicados en posiciones pares                                     & Lista                 & 3 \\
        Posiciones impares   & Lista             & Devuelve los elementos ubicados en posiciones impares                                   & Lista                 & 3 \\
        \hline
    \end{tabularx}
\end{table}

\begin{table}[H]
    \centering
    \caption{Nodos de Listas: Consulta y Búsqueda}
    \begin{tabularx}{\textwidth}{ l l X X c }
        \hline
        \textbf{Nodo}        & \textbf{Entradas} & \textbf{Proceso}                                                                  & \textbf{Salida}       & \textbf{Nivel de Abstracción} \\ \hline
        Contar               & Lista             & Calcula cuántos elementos tiene una lista guardada                                & Cantidad de elementos & 3 \\
        Primer elemento      & Lista             & Obtiene el primer elemento de la lista de entrada                                 & Elemento              & 1 \\
        Último elemento      & Lista             & Obtiene el último elemento de la lista de entrada                                 & Elemento              & 1 \\
        Elemento en posición & Lista, Posición   & Obtiene el elemento ubicado en la posición indicada de la lista (empezando en 1)  & Elemento              & 2 \\
        Buscar elemento      & Lista, Elemento   & Devuelve la primera posición donde aparece el elemento, o 0 si no aparece         & Posición              & 2 \\
        Lista vacía          & Lista             & Evalúa si la lista no contiene elementos                                          & Booleano              & 1 \\
        Mínimo               & Lista             & Obtiene el menor valor de la lista                                                & Elemento              & 2 \\
        Máximo               & Lista             & Obtiene el mayor valor de la lista                                                & Elemento              & 2 \\
        \hline
    \end{tabularx}
\end{table}

\begin{table}[H]
    \centering
    \caption{Nodos de Listas: Inserción y Eliminación}
    \begin{tabularx}{\textwidth}{ l l X X c }
        \hline
        \textbf{Nodo}     & \textbf{Entradas} & \textbf{Proceso}                                           & \textbf{Salida} & \textbf{Nivel de Abstracción} \\ \hline
        Añadir al inicio  & Lista, Elemento   & Inserta el elemento al inicio de la lista                  & Lista           & 1 \\
        Añadir al final   & Lista, Elemento   & Inserta el elemento al final de la lista                   & Lista           & 1 \\
        Quitar el primero & Lista             & Devuelve la lista de entrada eliminando su primer elemento & Lista           & 1 \\
        Quitar el último  & Lista             & Devuelve la lista de entrada eliminando su último elemento & Lista           & 1 \\
        \hline
    \end{tabularx}
\end{table}

\begin{table}[H]
    \centering
    \caption{Nodos de Listas: Agregación y Análisis}
    \begin{tabularx}{\textwidth}{ l l X X c }
        \hline
        \textbf{Nodo}   & \textbf{Entradas} & \textbf{Proceso}                                         & \textbf{Salida}      & \textbf{Nivel de Abstracción} \\ \hline
        Acumular        & Lista             & Suma todos los elementos de la lista                     & Total acumulado      & 2 \\
        Productorio     & Lista             & Multiplica todos los elementos de la lista               & Producto total       & 3 \\
        Frecuencia      & Lista             & Cuenta ocurrencias y devuelve pares (Elemento, Cantidad) & Lista de frecuencias & 3 \\
        \hline
    \end{tabularx}
\end{table}

\begin{table}[H]
    \centering
    \caption{Nodos de Lógica}
    \begin{tabularx}{\textwidth}{ l X X X c }
        \hline
        \textbf{Nodo}    & \textbf{Entradas}      & \textbf{Proceso}                                                                   & \textbf{Salida} & \textbf{Nivel de Abstracción} \\ \hline
        Comparar         & Valor A, Valor B       & Evalúa si A es mayor, menor, igual o diferente a B (según configuración).          & Booleano        & 2 \\
        Igual            & Valor A, Valor B       & Evalúa si A es igual a B.                                                          & Booleano        & 1-2 \\
        Distinto         & Valor A, Valor B       & Evalúa si A es distinto de B.                                                      & Booleano        & 1-2 \\
        Menor            & Valor A, Valor B       & Evalúa si A es menor que B.                                                        & Booleano        & 1-2 \\
        Menor o igual    & Valor A, Valor B       & Evalúa si A es menor o igual que B.                                                & Booleano        & 1-2 \\
        Mayor            & Valor A, Valor B       & Evalúa si A es mayor que B.                                                        & Booleano        & 1-2 \\
        Mayor o igual    & Valor A, Valor B       & Evalúa si A es mayor o igual que B.                                                & Booleano        & 1-2 \\
        Operación lógica & Booleano A, Booleano B & Aplica la compuerta seleccionada en configuración (AND, NAND, OR, NOR, XOR o NOT). & Booleano        & 2 \\
        AND              & Booleano A, Booleano B & Verifica si ambas entradas son Verdaderas.                                         & Booleano        & 1-2 \\
        NAND             & Booleano A, Booleano B & Negación de las salidas del nodo AND.                                              & Booleano        & 1-2 \\
        OR               & Booleano A, Booleano B & Verifica si al menos una entrada es Verdadera.                                     & Booleano        & 1-2 \\
        NOR              & Booleano A, Booleano B & Negación de las salidas del nodo OR.                                               & Booleano        & 1-2 \\
        XOR              & Booleano A, Booleano B & Verifica si solo una entrada es Verdadera.                                         & Booleano        & 1-2 \\
        NOT              & Booleano A             & Invierte el valor lógico de la entrada.                                            & Booleano        & 1-2 \\
        \hline
    \end{tabularx}
\end{table}

\begin{table}[H]
    \centering
    \caption{Nodos de Control de Flujo}
    \begin{tabularx}{\textwidth}{ l X X X c }
        \hline
        \textbf{Nodo}          & \textbf{Entradas}                                                                            & \textbf{Proceso}                                                                                                                                                                 & \textbf{Salida}                        & \textbf{Nivel de Abstracción} \\ \hline
        Desvío                 & Dato (Cualquiera), Condición (Booleano)                                                      & Si la condición es verdadera envía el dato por la salida superior; si es falsa, por la inferior.                                                                                 & Salida Sí (Dato), Salida No (Dato)     & 1 \\
        Relevo                 & Dato de activación, Dato a emitir                                                            & Usa una entrada para activar y emite siempre el valor conectado al puerto de relevo.                                                                                             & Dato relevado                          & 1 \\
        Selector               & Condición, Valor Verdadero, Valor Falso                                                      & Según la condición, emite el valor de la rama Verdadera o Falsa.                                                                                                                 & Valor seleccionado                     & 1 \\
        Switch                 & Dato (Cualquiera), Condición (Booleano)                                                      & Actúa como un desvío de tren: si la Condición es Verdadera, envía el Dato por la salida "Sí"; si es Falsa, por la salida "No".                                                   & Salida Sí (Dato), Salida No (Dato)     & 2 \\
        Condicional + Selector & Al condicional: Valores usados para las condiciones. \newline Al selector: Posibles salidas. & Evalúa las condiciones usando los valores de entrada del nodo condicional. Empareja cada condición con una salida del selector. Si una condición es verdadera, emite esa salida. & Número, texto o booleano seleccionado. & 3 \\
        \hline
    \end{tabularx}
\end{table}

\begin{table}[H]
    \centering
    \caption{Nodos de Distribución}
    \begin{tabularx}{\textwidth}{ l X X X c }
        \hline
        \textbf{Nodo}  & \textbf{Entradas}                     & \textbf{Proceso}                                                                                   & \textbf{Salida}                 & \textbf{Nivel de Abstracción} \\ \hline
        Replicar flujo & Flujo, Entradas para cada dispositivo & Define un flujo reutilizable y su configuración para ejecutarlo en varios dispositivos o destinos. & Resultados de todos los flujos  & 3 \\
        Distribuir     & Lista                                 & Distribuye los elementos de una lista entre múltiples salidas.                                     & Múltiples salidas con elementos & 2 \\
        Reunir         & Múltiples entradas                    & Recolecta valores de varias entradas y los agrupa en una lista.                                    & Lista con valores reunidos      & 2 \\
        \hline
    \end{tabularx}
\end{table}

\begin{table}[H]
    \centering
    \caption{Nodos de Micro:bit}
    \begin{tabularx}{\textwidth}{ l X X X c }
        \hline
        \textbf{Nodo}   & \textbf{Entradas}    & \textbf{Proceso}                                                              & \textbf{Salida}     & \textbf{Nivel de Abstracción} \\ \hline
        Leer botón      & Ninguna              & Espera por una pulsación de botón en la micro:bit y envía el valor detectado. & Texto (A, B o AB)   & 3 \\
        Temperatura     & Ninguna              & Lee la temperatura actual del dispositivo micro:bit.                          & Número (°C)         & 3 \\
        Leer pin        & Ninguna              & Lee el valor de un pin digital o analógico de la micro:bit.                   & Número              & 3 \\
        Escribir        & Dato (Cualquiera)    & Convierte la entrada a texto y lo escribe en las LED.                         & Sin salida de datos & 3 \\
        Emitir sonido   & Frecuencia, duración & Activa un sonido en micro:bit con la frecuencia y duración configuradas.      & Sin salida de datos & 3 \\
        Escribir en pin & Número (valor)       & Escribe un valor en un pin de salida de la micro:bit.                         & Sin salida de datos & 3 \\ \hline
    \end{tabularx}
\end{table}

\begin{table}[H]
    \centering
    \caption{Nodos de Abstracción}
    \begin{tabularx}{\textwidth}{ l X X X c }
        \hline
        \textbf{Nodo} & \textbf{Entradas}       & \textbf{Proceso}                                     & \textbf{Salida} & \textbf{Nivel de Abstracción} \\ \hline
        Subflujo      & Entradas personalizadas & Ejecuta un grupo de nodos encapsulados dentro de él. & Salida definida & 2 \\ \hline
    \end{tabularx}
\end{table}

\endgroup